% Methods: environment and reproducibility, the generated platform, the ten-gate
% design, the wired standards corpus, the pipeline and iteration design, and the
% verification floor. Dense facts live in left-aligned, body-width tabularx
% tables; long paths are given once near a table and abbreviated in cells.
% Figure 05 (standards wheel) is placed here and referenced with \autoref.

The methods are organised so the prose carries the argument and the tables carry
the dense facts. Each subsection names the repository record it draws from once,
then summarises the numbers in a left-aligned table set to the body measure and
refers to it rather than narrating every row. Across the section the same
discipline recurs: a reproducible model first, then the platform it represents,
then the gates that judge it, then the standards that justify the gate bounds,
then the pipeline that produces candidates, and finally the test floor that keeps
the whole apparatus honest.

\subsection{Environment and reproducibility}

The execution host was a managed ephemeral cloud container running Linux 6.18.5
on x86\_64, with CPython 3.11.15 as the single local interpreter and Claude Code
Opus 4.7 (1M context) Max as the autonomous runner. The dependency posture was
deliberately lean. The generated tree declares a core set of numpy, click,
pyyaml, and jsonschema, and these were installed alongside the ruff, yamllint,
and pytest tooling so the full check matrix could run. The heavy and
network-bound optional groups, the on-prem language-model backends anthropic,
ollama, and openai and the live Zenodo client, were intentionally left
uninstalled, so every module exercised its guarded offline path, which is the
honest continuous-integration behavior and the path a reviewer can reproduce
without credentials. Cross-version coverage on Python 3.10, 3.11, and 3.12 is
delegated to the integration matrix and is taken up in the Limitations section.

Reproducibility is the credibility floor, because a defensible assurance argument
begins with a model that regenerates exactly. Three committed artifacts, the
thousand-row positional sensor stream, the four-entrant comparison, and the sweep
index, were regenerated from scratch on the host and compared against the
committed versions; all three matched, the sensor stream byte-for-byte and the
other two as identical semantic reproductions. This byte-level determinism from a
fixed seed is precisely what NASA-STD-7009A and ASME V\&V 40-2018 \S3 expect
before any verification or validation claim can be called credible
\cite{nasastd7009a, asmevv40}. \autoref{tab:determinism} lists the regenerated
artifacts and the commands that produced them; all code was run from a scratch
directory with the generated tree on the import path, so the repository itself
stayed pristine throughout.

\begin{table}[H]
\caption{Determinism check, the reproducibility floor. Commands and artifact
paths are relative to \texttt{papers/VVUQ-02/codegen/}; all three regenerated
from seed 20260525.}
\label{tab:determinism}
\centering
\footnotesize
\begin{tabularx}{\textwidth}{L{3.2cm} Y L{3.5cm}}
\toprule
\textbf{Artifact (leaf)} & \textbf{Regeneration command} & \textbf{Result vs committed} \\
\midrule
\texttt{data/sample\_h2\_sensor.csv} & \texttt{python -m src.sensors.ingest\_h2} & byte-for-byte identical \\
\texttt{results/comparison.json} & \texttt{python -m src.llm.compare\_agent} & identical semantic match \\
\texttt{data/iterations/index.jsonl} & \texttt{python -m src.simulation.iterate --seed 20260525 --iterations 32} & identical \\
\bottomrule
\end{tabularx}
\end{table}

\subsection{The generated platform}

The generated tree models the hypothetical 2030 Unitree H2-Surgical 1.0, a
bipedal humanoid standing at the operating table whose base form factor traces to
the real Unitree H2 lineage and whose clinical-grade capability uplift is
explicitly flagged hypothetical and paper-only. The robot carries 71 total
degrees of freedom, distributed as fourteen across the two seven-axis arms, forty
across the two hands of five four-joint fingers, fifteen across the legs and
waist for stance, and two for the head. It commands its hands at a 10 kHz control
rate while sampling fingertip force at 100 kHz, and it places a fingertip to a
nominal accuracy near 0.05 mm under onboard vision. The platform performs the
60-second eight-phase Whipple on patient PAT-PDAC-0001, from the Kocher
exploration through the venous and uncinate dissections, transection and specimen
removal, the three reconstructive anastomoses, and closure. \autoref{tab:platform}
collects the morphology, timing, and force facts that the gates later check, and
\autoref{tab:config} maps the frozen configuration set that fixes them.

\begin{table}[H]
\caption{H2-Surgical 1.0 platform specification (hypothetical 2030), as fixed in
\texttt{papers/VVUQ-02/codegen/config/project.yaml}.}
\label{tab:platform}
\centering
\footnotesize
\begin{tabularx}{\textwidth}{L{5.2cm} Y}
\toprule
\textbf{Property} & \textbf{Value} \\
\midrule
Total degrees of freedom & 71 (14 arms, 40 hands, 15 stance, 2 head) \\
Arms and hands & 2 arms at 7 DOF; 2 hands at 20 DOF (5 fingers, 4 joints each) \\
Control and force rates & 10 kHz command; 100 kHz fingertip force \\
Fingertip placement accuracy & 0.05 mm RMS under onboard vision \\
Per-fingertip force caps & 2.5 N soft, 3.0 N hard \\
Bimanual cumulative force caps & 5.0 N soft, 6.0 N hard \\
Balance support-polygon margin floor & 8.0 mm \\
Safety latencies & 10 us proximity stop, 50 us joint park, 3 ms hand-arm e-stop \\
Procedure & pancreaticoduodenectomy, pylorus preserving, 8 phases, 3 anastomoses, 60 s \\
\bottomrule
\end{tabularx}
\end{table}

The source it runs from is organised into thirteen subpackages, namely
kinematics, sensors, perception, autonomy, hands, balance, safety, suturing,
vvuq, simulation, metrics, llm, and zenodo, each exercised in the execution
record and each mapped to the behavior or assurance role it plays. The
configuration set in \autoref{tab:config} is what makes those modules
deterministic: the kinematic chains, the hand and balance models, the perception
channels, the reference plan, the safety volumes, the anastomosis targets, and
the gate thresholds are all data, not code, so a re-generation from the same seed
yields the same platform.

\begin{table}[H]
\caption{Frozen configuration set; leaf names under the common parent
\texttt{papers/VVUQ-02/codegen/config/}.}
\label{tab:config}
\centering
\footnotesize
\begin{tabularx}{\textwidth}{L{4cm} Y}
\toprule
\textbf{Config file (leaf)} & \textbf{Defines} \\
\midrule
\texttt{project.yaml} & frozen scope: patient, seed 20260525, 32 iterations, phases, composite weights \\
\texttt{h2\_kinematics.yaml} & the two seven-axis arm chains plus head, waist, and leg chains \\
\texttt{hand\_model.yaml} & per-finger four-joint model, tendon coupling, and force caps \\
\texttt{balance\_model.yaml} & mass, centre of mass, support polygon, and ZMP margins \\
\texttt{perception\_model.yaml} & camera intrinsics and the NIR, ultrasound, and bile channels \\
\texttt{autonomy\_plan.yaml} & the eight-phase reference plan for concordance scoring \\
\texttt{safety\_zones.yaml} & the vessel no-fly keep-out volumes \\
\texttt{shared\_or\_actors.yaml} & human actor poses and keep-out spheres \\
\texttt{anastomosis\_targets.yaml} & the three ring-tension targets \\
\texttt{vvuq\_thresholds.yaml} & the ten gate threshold sets and the hard predicates \\
\texttt{standards\_map.yaml} & the machine-readable gate-to-standard binding \\
\bottomrule
\end{tabularx}
\end{table}

\subsection{The ten-gate assurance design}

The methodological core is the gate, and there are ten of them, one for each
named failure mode of a two-handed standing surgeon: bimanual hand-eye servoing,
dexterous finger force, whole-body balance, autonomous plan correctness,
instrument grasp and handover, the vascular no-fly hand, bimanual suturing and
anastomosis, perception scene understanding, shared-operating-room human
collision, and fault e-stop graceful degradation. Each gate is a fixed
composition of four steps applied in order. Verify checks that every structural
and policy predicate holds and requires a pass fraction of exactly 1.0, so a
single failed check blocks; verification is pass or fail, never a weighted
average. Validate compares the observed metric against an independent reference
committed under the reference data directory and requires agreement at or above a
per-gate bound within a maximum relative error, with a recorded human review.
Quantify computes the coefficient of variation across three or more seeded runs
and requires it to fall under a per-gate dispersion bound. Evaluate composes the
three and, on the three immediate-catastrophe gates, applies one extra hard
predicate that must also hold. The result is accept when every dimension passes,
block on the first failing dimension, or escalate when dispersion alone exceeds
its bound, which defaults the autonomy and fault paths to hand-back-to-human.

The bounds are not invented. They are set in proportion to model risk in the
manner ASME V\&V 40-2018 prescribes, so the three catastrophe gates carry the
tightest agreement, relative-error, and dispersion bounds and the extra hard
predicate, while the lower-risk behavior gates relax toward the defaults
\cite{asmevv40}. The hard predicates and the recorded human review are the risk
controls that ISO 14971:2019 expects for a high-severity hazard \cite{iso14971}.
The full per-gate threshold values are reported alongside the forest plot in the
Results section, in \autoref{tab:thresholds}, and are not reproduced here; the
point of this subsection is the composition and the risk-proportionate logic, not
the individual numbers.

\subsection{The wired external-standards corpus}

Credibility comes from binding each gate to published consensus standards already
used in real life, resolved at runtime rather than asserted in prose. The gate
registry loads a machine-readable manifest and attaches to each gate the
governing designations and the clause-derived limits it enforces, so every
decision the harness records carries its standards alongside its scores. The
input corpus holds fourteen consensus standards across model credibility,
robotic surgery, service-robot safety, software and risk, and autonomy, together
with two clinical baselines, the 2025 Dutch nationwide robotic
pancreatoduodenectomy cohort used as the human comparator and the Callery Fistula
Risk Score used to ground anastomosis quality. \autoref{fig:wheel} organises the
whole corpus into six domains radiating from the assurance hub, with the issuing
body on each segment and a regulatory relevancy callout naming what a real
deployment would require. The standards introduced here, and exercised across the
gates, are the ASME and NASA credibility standards \cite{asmevv40, nasastd7009a},
the robotic-surgery and medical-electrical standards \cite{iec8060127, iec60601},
the service-robot and industrial-robot safety standards \cite{isots15066,
iso13482, iso10218, iso9283}, the software life-cycle, risk, and control-system
standards \cite{iec62304, iso14971, iso13849}, the autonomy standards
\cite{ul4600, ieee7009}, and the United States Food and Drug Administration
computational-model credibility guidance \cite{fda2023credibility}. The reason to
wire them rather than cite them in passing is auditability: a reviewer or
regulator can follow any gate to its governing standard, which is what makes the
argument defensible and lets an inexpensive autonomous run stand as evidence.

\subsection{The control pipeline and the iteration design}

The candidates the gates judge come from a control pipeline that runs intent to
compile to act to score across six behavior groups. An on-prem language-model
step proposes phase-level intents, a deterministic policy compiles each intent
into a schema-conformant command of a Cartesian fingertip target, twenty
per-finger joint targets, and a grasp state, the behavior models act, and the
result feeds the six-component composite that the assurance layer then gates. On
this host the language-model step returned its offline fallback, the honest
continuous-integration path, and emitted the frozen reference plan, so the
phase-step concordance against that reference is 1.000 by construction. The
iteration design wraps the pipeline in a 32-iteration Latin hypercube at seed
20260525 over five free parameters, fingertip placement noise, finger force
noise, balance disturbance amplitude, perception occlusion fraction, and ring
tension drift, and the six-component composite reuses the frozen pancreatic
cancer weights of quality 0.30, time 0.20, cost 0.15, safety 0.15, patient
experience 0.05, and anastomosis quality 0.15. \autoref{tab:anchoring} reproduces
the binding of each behavior group to its gate and governing standards, so the
generation step is itself anchored, not just the assurance step that follows it.

\begin{table}[H]
\caption{External standards anchored at the pipeline stage; each behavior group
maps to its gate and governing standards (from
\texttt{papers/VVUQ-02/execution/02-pipeline/}).}
\label{tab:anchoring}
\centering
\footnotesize
\begin{tabularx}{\textwidth}{L{4.2cm} L{2.2cm} Y}
\toprule
\textbf{Behavior group} & \textbf{Gate} & \textbf{Governing standards} \\
\midrule
Autonomy intent and compile & VVUQ 04 & UL 4600 (2023), IEEE 7009-2024, IEC 62304 \\
Forward kinematics & VVUQ 01 & ISO 9283:1998, IEC 80601-2-77:2019 \\
Perception and fusion & VVUQ 08 & ASME V\&V 40-2018, IEC 62304 \\
Fingertip force and grasp & VVUQ 02, 05 & ISO/TS 15066:2016, IEC 80601-2-77:2019 \\
Whole-body balance & VVUQ 03 & ISO 13482:2014, IEC 60601-1 \\
Suturing and anastomosis & VVUQ 07 & IEC 80601-2-77:2019, Fistula Risk Score \\
\bottomrule
\end{tabularx}
\end{table}

\subsection{The verification floor}

Underneath the decision surface sits a shipped automated test suite that is the
verification floor in the sense of ASME V\&V 40-2018 \S8 and IEC 62304 software
verification \cite{iec62304}. The suite is 172 tests across fifteen modules, and
its weighting is itself evidence for the thesis: the assurance suite for the
ten-gate registry alone carries 64 of the 172 tests, more than a third of the
budget, and the safety module adds a further twelve, so the assurance and safety
modules dominate the behavior modules they judge. \autoref{tab:tests} gives the
per-module breakdown. The suite does not merely check that functions return; it
pins the safety semantics each standard requires, asserting that the catastrophe
gates produce a hard stop, an e-stop, or a hand-back at the boundary, that the
balance gate holds its zero-moment-point margin, that the autonomy gate scores
plan concordance, and that the standards-corpus tests resolve every binding to a
published designation. The suite runs in about 1.4 seconds on the host, and it
must be green before any sweep or decision-surface run is trusted.

\begin{table}[H]
\caption{Per-module automated test suite, the verification floor; module leaf
names under \texttt{papers/VVUQ-02/codegen/tests/}. Total 172 tests, 0 skipped.}
\label{tab:tests}
\centering
\footnotesize
\begin{tabularx}{\textwidth}{L{4.2cm} L{1.4cm} Y}
\toprule
\textbf{Test module (leaf)} & \textbf{Tests} & \textbf{Subject} \\
\midrule
\texttt{test\_vvuq\_gates.py} & 64 & the ten-gate registry, thresholds, and decision surface \\
\texttt{test\_safety.py} & 12 & e-stop fault monitor, human-collision FSM, vessel no-fly gate \\
\texttt{test\_suturing.py} & 11 & bimanual suture, ring-tension restore, in-band metrics \\
\texttt{test\_hands.py} & 11 & fingertip force tracking, bimanual caps, grasp, handover \\
\texttt{test\_vvuq\_framework.py} & 9 & verify, validate, quantify, and gate composition \\
\texttt{test\_kinematics.py} & 9 & forward kinematics, joint-limit clamp, distance \\
\texttt{test\_perception.py} & 8 & scene segmentation, Dice, and channel fusion \\
\texttt{test\_balance.py} & 8 & ZMP computation, support-polygon margin, recovery \\
\texttt{test\_autonomy.py} & 8 & intent proposal, deterministic compile, concordance \\
\texttt{test\_sensors.py} & 7 & per-tick per-hand sensor synthesis and serialization \\
\texttt{test\_llm.py} & 7 & the four-entrant tournament agent and caveats \\
\texttt{test\_simulation.py} & 6 & the 32-iteration Latin hypercube sweep \\
\texttt{test\_metrics.py} & 5 & the six-component composite and the gating overlay \\
\texttt{test\_zenodo.py} & 4 & L0 pointer JSON and the cross-iteration manifest \\
\texttt{test\_standards\_corpus.py} & 3 & the wired external-standards manifest binding \\
\bottomrule
\end{tabularx}
\end{table}

\begin{figure}[H]
   \centering
   \IfFileExists{Images/05-clinical-regulatory-standards-wheel.png}{%
     \hspace{-0.95cm}%
     \includegraphics[width=1.1\linewidth, trim={0 0 0 0}, clip]{Images/05-clinical-regulatory-standards-wheel.png}%
     % \vspace{-0.25cm}
   }{%
     \setlength{\fboxsep}{10pt}%
     \fbox{\parbox[c][6cm][c]{0.9\linewidth}{\centering\itshape Figure placeholder.
     The author adds \texttt{Images/05-clinical-regulatory-standards-wheel.png},
     rendered by
     \texttt{papers/VVUQ-02/imagegen/05-clinical-regulatory-standards-wheel/05-clinical-regulatory-standards-wheel.py}.}}%
   }
    \caption{Clinical and regulatory standards wheel for the wired VVUQ-02 input corpus of fourteen consensus standards and two clinical baselines across six domains.}
    \label{fig:wheel}
\end{figure}
